Language Identification in the (Pretraining) Data Processing
Pipeline – State of the Art and Open Problems

The first step in training a broad-coverage multilingual
language model is to run a pipeline that produces its
pretraining data.  We address the problem of language
identification (LID) in such a pipeline.  Many challenges in
this setting are inherent in the complexities of
multilinguality: thousands of classes, many of which are hard
to distinguish in pretraining data (e.g., Malay vs
Indonesian).  Other challenges are due to a severe mismatch
between the training data available for LID and the
pretraining data. This mismatch also applies to evaluation:
it is often not feasible to directly evaluate the
multilingual quality of the trained language model (the
quantity we care about). LID measures like F1 and false
positive rate on held-out validation data may not be
strongly correlated with LM quality. Finally, I will present
GlotLID, our LID model for creating multilingual pretraining
data, including our design decisions for addressing some of
the challenges of the problem. This includes separate
"language class" labels for different noise categories,
introducing composite labels  (ISO-639 and script),
reducing the number of confusable labels and optimizing the
speed-accuracy tradeoff for pretraining data generation.
